\documentclass[11pt,a4paper]{article}

\usepackage[utf8]{inputenc}
\usepackage[T1]{fontenc}
\usepackage[french]{babel}
\usepackage{amsmath,amssymb}
\usepackage{graphicx}
\usepackage[dvipsnames]{xcolor}
\usepackage{float}
\usepackage{booktabs}
\usepackage[margin=1.5cm,top=2cm]{geometry}
\usepackage{hyperref}

\graphicspath{{../test_low_mach_compar_resume/}}
\newcommand{\amndir}{../test_low_mach_compar_resume/}
\newcommand{\sch}{WIP}
\newcommand{\schpath}{WIP}

\title{Schéma \texttt{\sch} --- suite bas Mach / hypersonique}
\author{Vincent Delmas}
\date{\today}

\begin{document}
\maketitle

\noindent
Résultats de \texttt{subfvns} pour le schéma \texttt{\sch}, sur la suite complète
allant du tourbillon de Gresho au flux de chaleur pariétal à Mach~27.

\medskip
\noindent\textbf{Provenance des résultats.} Tous les cas de tous les schémas comparés
(\texttt{multi\_point}, \texttt{three\_wave}, \texttt{multi\_point\_pressure},
\texttt{WIP}, \texttt{WIP2\_NOLM}) proviennent d'une \emph{seule} campagne lancée avec
un \emph{seul} binaire, dans \texttt{test\_low\_mach\_compar\_resume/}. Les chiffres des
cinq rapports sont donc directement comparables entre eux, ce qui n'était pas le cas
des campagnes précédentes, étalées sur plusieurs reconstructions du code.

\medskip
\noindent\textbf{Critères évalués.}
\begin{enumerate}\itemsep2pt
  \item préservation du champ de vitesse sur le tourbillon de Gresho (quad et tri) ;
  \item convergence de la fluctuation de densité en fonction du Mach --- au moins
        $\mathrm{Ma}^1$, idéalement $\mathrm{Ma}^2$ ;
  \item robustesse sur Sedov, maillages triangulaire \emph{et} hexaédrique ;
  \item flux de chaleur pariétal sur le demi-cylindre à Mach~27, quad et tri.
\end{enumerate}

\medskip
\noindent\textbf{Deux réserves de lecture.} L'échelle absolue du nombre de Stanton
provient de la formule déjà présente dans \texttt{plot\_st.gnu} et n'a pas été
re-dérivée ici : les classements entre schémas sont fiables, la valeur absolue
mériterait une vérification. Par ailleurs les maillages triangulaires sont produits
par la version de \texttt{gmsh} présente sur le cluster ; ils diffèrent de ceux
obtenus localement, de sorte que seules les comparaisons \emph{à maillage identique}
ont un sens --- ce qui est le cas ici, tous les schémas partageant les mêmes maillages.

\clearpage
\section*{Résultats}

\newcommand{\minifig}[2]{%
  \begin{minipage}[t]{0.48\linewidth}\centering
    \IfFileExists{\amndir#1}%
      {\includegraphics[width=\linewidth]{#1}}%
      {\fbox{\parbox[c][0.3\linewidth]{0.8\linewidth}{\centering\small non disponible}}}%
    \\[2pt]{\small #2}
  \end{minipage}\hfill}

\subsection*{Gresho -- quad}
\begin{figure}[H]\centering
  \minifig{gresho_quad/outputs/\schpath/images/gresho_quad_velocity.png}{champ de vitesse}
  \minifig{gresho_quad/outputs/\schpath/images/velocity_profile.pdf}{profil radial}
\end{figure}

\subsection*{Gresho -- tri}
\begin{figure}[H]\centering
  \minifig{gresho_tri/outputs/\schpath/images/gresho_tri_velocity.png}{champ de vitesse}
  \minifig{gresho_tri/outputs/\schpath/images/velocity_profile.pdf}{profil radial}
\end{figure}

\subsection*{Gresho -- Voronoï}
Maillage polygonal irrégulier de 2500 cellules (même nombre que le quad
$50\times50$, donc à résolution égale), 4 à 8 côtés, engendré par
\texttt{tools/voronoi/voronoi\_box.py}. Les germes étant construits par
réflexion à travers les deux axes médians pour que la périodicité soit exacte au
niveau machine, ce maillage est \emph{miroir-symétrique} par rapport à $x=0$ et
$y=0$ : ce n'est pas un Voronoï générique, et la couture peut transparaître dans
les champs très dissipés.
\begin{figure}[H]\centering
  \minifig{gresho_voro/outputs/\schpath/images/gresho_voro_velocity.png}{champ de vitesse}
  \minifig{gresho_voro/outputs/\schpath/images/velocity_profile.pdf}{profil radial}
\end{figure}
\noindent\footnotesize\emph{Attention à la lecture :} l'échelle de couleur est
ajustée automatiquement sur chaque image. Pour un schéma qui a détruit le
tourbillon, elle amplifie un bruit résiduel d'énergie quasi nulle, qui ne doit
pas être confondu avec un champ de vitesse de même amplitude que celui d'un
schéma qui l'a préservé.\normalsize

\subsection*{Convergence bas Mach de la fluctuation de densité}
\begin{figure}[H]\centering
  \minifig{convergence_gresho_quad/outputs/\schpath/images/convergence.pdf}{quad}
  \minifig{convergence_gresho_tri/outputs/\schpath/images/convergence.pdf}{tri}
\end{figure}
\begin{figure}[H]\centering
  \minifig{convergence_gresho_voro/outputs/\schpath/images/convergence.pdf}{Voronoï}
\end{figure}

\clearpage
\subsection*{Sedov -- maillage triangulaire}
\begin{figure}[H]\centering
  \minifig{sedov_tri/outputs/\schpath/images/sedov_tri_density.png}{densité}
  \minifig{sedov_tri/outputs/\schpath/density_profile.png}{profil radial}
\end{figure}

\subsection*{Sedov -- maillage hexaédrique}
\begin{figure}[H]\centering
  \minifig{sedov_hex/outputs/\schpath/images/sedov_hex_density.png}{densité}
  \minifig{sedov_hex/outputs/\schpath/density_profile.png}{profil radial}
\end{figure}

\clearpage
\subsection*{Demi-cylindre Navier-Stokes -- quad}
\begin{figure}[H]\centering
  \minifig{half_cylinder_quad_ns/outputs/\schpath/images/output_-1_Pressure.png}{pression}
  \minifig{half_cylinder_quad_ns/outputs/\schpath/images/output_-1_Temperature.png}{température}
\end{figure}
\begin{figure}[H]\centering
  \minifig{half_cylinder_quad_ns/outputs/\schpath/plot_cp.pdf}{$C_p$}
  \minifig{half_cylinder_quad_ns/outputs/\schpath/plot_st.pdf}{flux de chaleur (Stanton)}
\end{figure}

\subsection*{Demi-cylindre Navier-Stokes -- maillage hybride quad + Voronoï}
Couche limite structurée en quadrangles sur $r\in[1;1{,}04]$ puis lointain en
Voronoï, conformes à l'interface, 9103 mailles --- la transposition directe du
maillage \texttt{tri}, qui est lui-même un hybride quad + triangles de 8326
mailles. C'est ce maillage qui \emph{discrimine} : sur quad, presque tous les
schémas se superposent à la référence LAURA et on ne les distingue pas.
\begin{figure}[H]\centering
  \minifig{half_cylinder_voro_ns/outputs/\schpath/images/output_-1_Pressure.png}{pression}
  \minifig{half_cylinder_voro_ns/outputs/\schpath/images/output_-1_Temperature.png}{température}
\end{figure}
\begin{figure}[H]\centering
  \minifig{half_cylinder_voro_ns/outputs/\schpath/plot_cp.pdf}{$C_p$}
  \minifig{half_cylinder_voro_ns/outputs/\schpath/plot_st.pdf}{flux de chaleur (Stanton)}
\end{figure}

\subsection*{Demi-cylindre Navier-Stokes -- tri}
\begin{figure}[H]\centering
  \minifig{half_cylinder_tri_ns/outputs/\schpath/images/output_-1_Pressure.png}{pression}
  \minifig{half_cylinder_tri_ns/outputs/\schpath/images/output_-1_Temperature.png}{température}
\end{figure}
\begin{figure}[H]\centering
  \minifig{half_cylinder_tri_ns/outputs/\schpath/plot_cp.pdf}{$C_p$}
  \minifig{half_cylinder_tri_ns/outputs/\schpath/plot_st.pdf}{flux de chaleur (Stanton)}
\end{figure}

\end{document}
